\section{Evaluating Applicability (D4)}\label{sec:evaluation}
\begin{table*}[htb]
    \centering
    \footnotesize
    \caption{Baseline evaluation of each risk in \method on \celeba, reproducing prior attacks to validate the implementations. $\modelstd$, $\modelpois$, and $\modelstol$ are the baseline, poisoned, and stolen models. $^{*}$Membership inference uses an intentionally overfit ResNet-18.}
    \label{tab:attack_results}
    \begin{tabularx}{\textwidth}{l l X X X X}
        \toprule
        \multirow{2}{*}{\textbf{Attack}} & \textbf{Model Architecture} & \textbf{VGG11} & \textbf{VGG13} & \textbf{VGG16} & \textbf{VGG19} \\
            & $Acc_{te}$ & 91.21~$\pm$~0.05 & 91.46~$\pm$~0.04 & 91.38~$\pm$~0.03 & 91.41~$\pm$~0.03 \\
        \midrule

        \ref{evasion}~(Evasion) & $Acc_{rob}$ & 10.08~$\pm$~0.13 & 8.60~$\pm$~0.05 & 8.65~$\pm$~0.02 & 8.67~$\pm$~0.04 \\
        \midrule

        \multirow{3}{*}{\ref{poison}~(Poisoning)}
            &  $Acc_{pois}$ ($\modelstd$) & 8.13~$\pm$~0.30 & 8.27~$\pm$~0.34 & 8.17~$\pm$~0.19 & 7.57~$\pm$~0.31 \\
            & $Acc_{te}$ ($\modelpois$) & 91.16~$\pm$~0.03 & 91.48~$\pm$~0.04 & 91.39~$\pm$~0.03 & 91.37~$\pm$~0.08 \\
            & $Acc_{pois}$ ($\modelpois$) & 95.70~$\pm$~0.21 & 96.70~$\pm$~0.10 & 96.62~$\pm$~0.18 & 96.75~$\pm$~0.18 \\
        \midrule

        \multirow{3}{*}{\ref{modelext}~(Unauthorized Model Ownership)}
            &  $Acc_{te}$ ($\modelstol$) & 91.02~$\pm$~0.09 & 91.32~$\pm$~0.07 & 91.42~$\pm$~0.08 & 91.39~$\pm$~0.06 \\
            & $Fid$ ($\modelstol$) & 94.64~$\pm$~0.16 & 94.75~$\pm$~0.09 & 94.74~$\pm$~0.11 & 94.84~$\pm$~0.08 \\
            & $Fid_{cor}$ ($\modelstol$) & 93.14~$\pm$~0.06 & 93.49~$\pm$~0.03 & 93.51~$\pm$~0.02 & 93.41~$\pm$~0.04 \\
        \midrule

        \multirow{8}{*}{\ref{meminf}~(Membership Inference)$^{*}$}
            & $Acc_{tr}$ & 94.20~$\pm$~1.33 & - & - & - \\
            & $Acc_{te}$ & 76.81~$\pm$~0.48 & - & - & - \\
            & Offline $Acc_{bal}$ & 50.41~$\pm$~00.13 & - & - & - \\
            & Offline $AUC$ & 0.49~$\pm$~0.00 & - & - & - \\
            & Offline TPR@1\%FPR & 0.83~$\pm$~0.11 & - & - & - \\
            & Online $Acc_{bal}$ & 60.67~$\pm$~00.70 & - & - & - \\
            & Online $AUC$ & 0.64~$\pm$~0.01 & - & - & - \\
            & Online TPR@1\%FPR & 2.70~$\pm$~0.25 & - & - & - \\
        \midrule

        \multirow{2}{*}{\ref{attinf}~(Attribute Inference)}
            & $Acc_{att}$ & 57.17~$\pm$~0.15 & 57.12~$\pm$~0.12 & 57.27~$\pm$~0.11 & 56.90~$\pm$~0.13 \\
            & $AUC$ & 0.59~$\pm$~0.00 & 0.59~$\pm$~0.00 & 0.59~$\pm$~0.00 & 0.58~$\pm$~0.00 \\
        \midrule

        \multirow{3}{*}{\ref{datarecon}~(Data Reconstruction)}
            & $MSE_{avg}$ & 0.20~$\pm$~0.00 & 0.20~$\pm$~0.00 & 0.20~$\pm$~0.00 & 0.20~$\pm$~0.000 \\
            & $MSE_0$ & 0.19~$\pm$~0.00 & 0.19~$\pm$~0.00 & 0.19~$\pm$~0.00 & 0.19~$\pm$~0.00 \\
            & $MSE_1$ & 0.20~$\pm$~0.00 & 0.20~$\pm$~0.00 & 0.20~$\pm$~0.00 & 0.20~$\pm$~0.00 \\
        \bottomrule
    \end{tabularx}
\end{table*}


To evaluate the applicability requirement (D4), we begin with baseline experiments on the \celeba dataset.
These baselines validate the correctness of the implementations and establish reference points for future comparisons.
We train VGG11/13/16/19 models for 100 epochs as our main architectures using cross-entropy loss with the Adam optimizer unless otherwise specified.
Each risk is then evaluated with its corresponding attack modules from Table~\ref{tab:modules}, and all experiments are repeated ten times with mean and standard error reported.
We focus on \celeba since it combines image data with sensitive attributes, allowing us to assess multiple risks using one dataset.
After presenting these baselines, we study three unintended interactions that prior work has left unexplored or only partially explained.

\subsection{Baseline Results}
Baseline results are shown in Table~\ref{tab:attack_results}.
We evaluate each algorithm using \method{'s} metrics and compare to the original papers when possible.
Overall, our reproductions closely match prior work, with exceptions noted below.

\noindent\textbf{\ref{evasion} (Evasion).}
Prior work evaluates PGD on \cifar and reports robust accuracy ($Acc_{rob}$) of $0.8$--$3.5\%$ for a ResNet model~\cite{madry2018DeepLearningModels}.
With a comparable setup using \cifar, we obtained $Acc_{rob}=3.8\%$, consistent with these results.
On \celeba, however, $Acc_{rob}$ is higher ($8.6$--$10\%$), making the attack less effective.
We suspect this difference arises because \celeba ($\approx202{,}000$ images) is much larger than \cifar ($60{,}000$ images), and larger datasets typically yield more robust models~\cite{schmidt2018AdversariallyRobustGeneralization}.

\noindent\textbf{\ref{poison} (Poisoning).}
The original work evaluated BadNets on MNIST and traffic signs~\cite{gu2019BadNetsEvaluatingBackdooring}.
Since MNIST is now considered too simple for robust evaluation, we reproduced it only for reference.
Using a two-layer CNN, we obtained a poisoned data accuracy ($Acc_{pois}$) of $99\%$.

On \celeba, the $\modelpois$ achieves $\approx96\%$ $Acc_{pois}$, demonstrating that the backdoor was successfully learned and the attack is effective.
In contrast, $\modelstd$ reaches only $\approx8\%$, since it never learned the trigger--target association.

\noindent\textbf{\ref{modelext} (Unauthorized Model Ownership).}
Prior work evaluated model extraction on \celeba with VGG19, combining three attributes into an 8-class classification problem~\cite{liu2022MLDoctorHolisticRisk}.
They reported a best fidelity of $\approx81\%$ between $\modelstol$ and $\modelstd$.

In contrast, \method achieves $\approx94\%$ fidelity on \celeba for the binary task, likely because the victim models in prior work were more overfit (Table 2 in~\cite{liu2022MLDoctorHolisticRisk}).
We also report \emph{correct fidelity} ($Fid_{cor}$) for completeness, which measures agreement conditioned on correctness, i.e., when models agree, how often are they also correct.

\noindent\textbf{\ref{meminf} (Membership Inference).}
To evaluate~\ref{meminf}, \method reports TPR at a fixed FPR of $1\%$, as recommended in prior work~\cite{carlini2022MembershipInferenceAttacks}.
We also report balanced accuracy (best threshold from the ROC curve) and ROC AUC.
Prior work evaluated this attack on \cifar with a ResNet model, using 256 shadow models, and achieved a TPR of $2.2\%$ at $0.001\%$ FPR, $8.4\%$ TPR at $0.1\%$ FPR, and a balanced accuracy of $63.8\%$~\cite{carlini2022MembershipInferenceAttacks}.

To reproduce comparable results, we trained an overfit ResNet-18 model, as shown in Table~\ref{tab:attack_results}.
In contrast, a well-trained \celeba model is not susceptible to this attack.
This is because membership inference is more successful against overfit models~\cite{carlini2022MembershipInferenceAttacks}.

\noindent\textbf{\ref{attinf} (Attribute Inference).}
Performance is measured using balanced accuracy ($Acc_{bal}$) and ROC AUC.
Prior work reported $53$--$62\%$ $Acc_{bal}$ across different datasets~\cite{aalmoes2024AlignmentGroupFairness}.
On \celeba, we obtain $\approx57\%$ accuracy (Table~\ref{tab:attack_results}).
In addition, our preliminary tests on \census achieve $\approx65\%$ (Race) and $\approx59\%$ (Sex), closely matching the original results of $\approx56\%$ and $\approx58\%$.
Overall, this suggests our results are consistent and reasonable.

\noindent\textbf{\ref{datarecon} (Data Reconstruction).}
Prior work applied the same reconstruction algorithm using a similar architecture on \celeba but with an 8-class setup (vs.\ our binary task) and reported a lower MSE ($\approx0.1$ vs. $\approx0.2$)~\cite{liu2022MLDoctorHolisticRisk}.
Because the attack recovers an \emph{average} data record per class, having more classes reduces intra-class variability and typically yields lower MSE, explaining the discrepancy.

% \input{tables/tab_sumInteractions}

\begin{table*}[!htb]
    \setlength\tabcolsep{2.5pt}
    \footnotesize
    \centering
    \caption{Interaction between~\ref{advtr} Adversarial Training and~\ref{attinf} Attribute Inference on \census and \lfw. Adversarial training does not consistently change susceptibility to~\ref{attinf}.}
    \label{tab:attinf_advtr}
    \begin{tabular}{ c c c c c c c c }
        \bottomrule

        \toprule
        \multicolumn{8}{c}{\textbf{\census}}\\
            & $Acc_{te}$ & $Acc_{rob}$ ($\modelstd$) & $Acc_{rob}$ ($\modeldef$) & $Acc_{att}$ (\race) & $AUC$ (\race) & $Acc_{att}$ (\sex) & $AUC$ (\sex) \\
        \midrule
          Baseline ($\modelstd$) & 81.99~$\pm$~0.57 & - & - & 58.58~$\pm$~0.88 & 0.61~$\pm$~0.01 & 65.06~$\pm$~1.21 & 0.69~$\pm$~0.01 \\
          $\epsilon_{rob}$ = 0.01 & 82.53~$\pm$~0.44 & 73.87~$\pm$~1.18 & 81.15~$\pm$~0.68 & 59.64~$\pm$~0.53 & 0.62~$\pm$~0.01 & 65.88~$\pm$~0.80 & 0.70~$\pm$~0.01 \\
          $\epsilon_{rob}$ = 0.03 & 83.00~$\pm$~0.26 & 44.25~$\pm$~4.27 & 80.10~$\pm$~0.79 & 60.23~$\pm$~0.63 & 0.63~$\pm$~0.01 & 66.23~$\pm$~0.58 & 0.70~$\pm$~0.00 \\
          $\epsilon_{rob}$ = 0.06 & 83.40~$\pm$~0.36 & 22.70~$\pm$~1.95 & 80.03~$\pm$~0.64 & 60.28~$\pm$~0.31 & 0.63~$\pm$~0.00 & 66.78~$\pm$~1.09 & 0.71~$\pm$~0.01 \\
          $\epsilon_{rob}$ = 0.1 & 83.70~$\pm$~0.24 & 20.89~$\pm$~0.88 & 76.03~$\pm$~2.58 & 60.60~$\pm$~0.59 & 0.64~$\pm$~0.01 & 67.44~$\pm$~1.07 & 0.72~$\pm$~0.01 \\
        \midrule
        \multicolumn{8}{c}{\textbf{\lfw}}\\
            & $Acc_{te}$ & $Acc_{rob}$ ($\modelstd$) & $Acc_{rob}$ ($\modeldef$) & $Acc_{att}$ (\race) & $AUC$ (\race) & $Acc_{att}$ (\sex) & $AUC$ (\sex) \\
        \midrule
          Baseline ($\modelstd$) & 82.84~$\pm$~0.56 & - & - & 66.80~$\pm$~4.07 & 0.73~$\pm$~0.05 & 76.59~$\pm$~1.00 & 0.84~$\pm$~0.01 \\
          $\epsilon_{rob}$ = 0.01 & 81.36~$\pm$~0.67 & 54.31~$\pm$~1.48 & 75.09~$\pm$~0.95 & 68.41~$\pm$~6.43 & 0.74~$\pm$~0.08 & 76.16~$\pm$~0.74 & 0.84~$\pm$~0.01 \\
          $\epsilon_{rob}$ = 0.03 & 79.57~$\pm$~0.59 & 18.66~$\pm$~0.48 & 50.22~$\pm$~2.39 & 68.64~$\pm$~3.65 & 0.76~$\pm$~0.05 & 75.22~$\pm$~0.91 & 0.82~$\pm$~0.02 \\
          $\epsilon_{rob}$ = 0.06 & 78.47~$\pm$~0.66 & 17.20~$\pm$~0.55 & 33.28~$\pm$~2.58 & 72.37~$\pm$~2.79 & 0.79~$\pm$~0.04 & 75.06~$\pm$~1.30 & 0.82~$\pm$~0.01 \\
          $\epsilon_{rob}$ = 0.1 & 81.29~$\pm$~0.92 & 17.15~$\pm$~0.57 & 19.25~$\pm$~0.85 & 67.79~$\pm$~5.38 & 0.74~$\pm$~0.05 & 75.78~$\pm$~1.42 & 0.82~$\pm$~0.01 \\
        \bottomrule

        \toprule
    \end{tabular}
\end{table*}


\subsection{Unintended Interaction}
We now use \method to identify the nature of different unexplored unintended interactions. We focus on
\begin{enumerate*}[label=\roman*),itemjoin={,\xspace}] % chktex 9 % chktex 10
\item \emph{unexplored interactions} and \item \emph{partially explored interactions} where the nature of the interaction is known, but the underlying factors that influence it are unclear.
\end{enumerate*}
% We summarize both unexplored and partially explored interactions in Table~\ref{tab:sumInteractions}.

We give three examples of unintended interactions below.

% Adversarial Training and Attribute Inference
\noindent\textbf{\ref{advtr} (Adversarial Training) and~\ref{attinf} (Attribute Inference).}
We use \census and \lfw to study how adversarial training (\ref{advtr}) influences attribute inference (\ref{attinf}), summarized in Table~\ref{tab:attinf_advtr}.
In this experiment, we reserve 50\% of the training data for the adversary, reducing $Acc_{te}$ slightly.
Unless otherwise stated, we regard changes smaller than 5 p.p. or within one standard error as not significant, reflecting normal variation across runs.

On \census, $Acc_{bal}$ for \race rises only slightly, from $\approx58\%$ at baseline to $\approx60\%$ at $\epsilon_{rob}=0.01$, after which it remains stable.
For the \sex attribute, $Acc_{bal}$ grows gradually from $\approx65\%$ to $\approx67\%$ at $\epsilon_{rob}=0.1$, with the $AUC$ following the same pattern.
However, in both attributes, the overall change is only about 2 p.p., indicating that the effect of adversarial training on \census is negligible.


For \lfw, the change in $Acc_{bal}$ is more significant.
The baseline $Acc_{bal}$ is $\approx66\%$ and increases to $\approx72\%$ for a $\epsilon_{rob} = 0.06$, and $AUC$ follows a similar trend.
Similar to the results for \census, the metrics for the attribute \sex have the opposite trend to the metrics for the attribute \race.
However, it should be noted that the adversarial training in \lfw was not as effective as it was for \census.
At $\epsilon_{rob} = 0.06$, $Acc_{rob}$ for $\modeldef$ is only $\approx33\%$.
Although $Acc_{rob}$ for $\modeldef$ is higher than $\modelstd$, it would be a stretch to call this a robust model.

Based on these experiments, changes in $Acc_{bal}$ and $AUC$ for a robust model fall well within the 5 p.p. range, indicating that~\ref{advtr} does not have a significant impact on~\ref{attinf}.
However, because $\modeldef$ is not strongly robust for \lfw and the results are inconsistent across attributes and datasets, we cannot conclusively determine whether~\ref{advtr} affects~\ref{attinf}.

\begin{table*}[htb]
    \setlength\tabcolsep{2.5pt}
    \centering
    \footnotesize
    \caption{Interaction between~\ref{advtr} Adversarial Training and~\ref{modelext} Unauthorized Model Ownership. $Fid$ and $Fid_{cor}$ are fidelity and correct fidelity between $\modeldef$ and the stolen model $\modelstol$. Adversarial training raises fidelity but lowers correct fidelity on harder datasets.}

    \label{tab:advtr_modelext}
    \begin{tabular}{ c c c c c c c }
        \bottomrule

        \toprule
        \multicolumn{7}{c}{\textbf{\census}}\\
            & $Acc_{te}$ ($\modeldef$) & $Acc_{rob}$ ($\modelstd$) & $Acc_{rob}$ ($\modeldef$) & $\modelstol$: $Acc_{te}$ & $\modelstol$: $Fid$ & $\modelstol$: $Fid_{cor}$ \\
        \midrule
        Baseline ($\modelstd$) & 81.95~$\pm$~0.62 & - & - & - & - & - \\
        $\epsilon_{rob}$ = 0.01 & 82.27~$\pm$~1.59 & 73.25~$\pm$~1.15 & 81.43~$\pm$~1.84 & 82.74~$\pm$~1.69 & 95.30~$\pm$~0.98 & 84.10~$\pm$~1.47 \\
        $\epsilon_{rob}$ = 0.03 & 82.88~$\pm$~0.74 & 42.40~$\pm$~5.58 & 80.33~$\pm$~1.69 & 83.30~$\pm$~0.62 & 96.44~$\pm$~0.74 & 84.31~$\pm$~0.51 \\
        $\epsilon_{rob}$ = 0.05 & 83.02~$\pm$~0.89 & 26.02~$\pm$~4.46 & 78.99~$\pm$~2.91 & 83.46~$\pm$~0.88 & 96.95~$\pm$~0.60 & 84.28~$\pm$~0.76 \\
        $\epsilon_{rob}$ = 0.10 & 84.06~$\pm$~0.15 & 20.32~$\pm$~1.65 & 79.23~$\pm$~2.45 & 84.38~$\pm$~0.15 & 98.34~$\pm$~0.38 & 84.79~$\pm$~0.21 \\
        \midrule
        \multicolumn{7}{c}{\textbf{\fmnist}}\\
            & $Acc_{te}$ ($\modeldef$) & $Acc_{rob}$ ($\modelstd$) & $Acc_{rob}$ ($\modeldef$) & $\modelstol$: $Acc_{te}$ & $\modelstol$: $Fid$ & $\modelstol$: $Fid_{cor}$ \\
        \midrule
        Baseline ($\modelstd$) & 85.67~$\pm$~0.31 & - & - & - & - & - \\
        $\epsilon_{rob}$ = 0.01 & 84.81~$\pm$~0.69 & 77.38~$\pm$~0.73 & 80.63~$\pm$~0.69 & 84.45~$\pm$~0.60 & 92.42~$\pm$~0.45 & 88.27~$\pm$~0.59 \\
        $\epsilon_{rob}$ = 0.03 & 84.92~$\pm$~0.43 & 42.54~$\pm$~1.31 & 74.59~$\pm$~0.72 & 84.32~$\pm$~0.37 & 94.76~$\pm$~0.34 & 87.07~$\pm$~0.36 \\
        $\epsilon_{rob}$ = 0.05 & 83.86~$\pm$~0.58 & 19.59~$\pm$~0.91 & 69.08~$\pm$~1.46 & 83.26~$\pm$~0.43 & 95.78~$\pm$~0.23 & 85.39~$\pm$~0.38 \\
        $\epsilon_{rob}$ = 0.10 & 80.58~$\pm$~0.63 & 7.97~$\pm$~0.33 & 58.55~$\pm$~2.10 & 80.30~$\pm$~0.51 & 96.16~$\pm$~0.37 & 82.15~$\pm$~0.46 \\
        \midrule
        \multicolumn{7}{c}{\textbf{\lfw}}\\
            & $Acc_{te}$ ($\modeldef$) & $Acc_{rob}$ ($\modelstd$) & $Acc_{rob}$ ($\modeldef$) & $\modelstol$: $Acc_{te}$ & $\modelstol$: $Fid$ & $\modelstol$: $Fid_{cor}$ \\
        \midrule
        Baseline ($\modelstd$) & 82.89~$\pm$~0.63 & - & - & - & - & - \\
        $\epsilon_{rob}$ = 0.01 & 81.04~$\pm$~1.22 & 52.51~$\pm$~1.79 & 74.80~$\pm$~2.48 & 78.73~$\pm$~3.48 & 87.91~$\pm$~4.13 & 84.00~$\pm$~1.34 \\
        $\epsilon_{rob}$ = 0.03 & 77.83~$\pm$~1.62 & 18.03~$\pm$~0.65 & 55.08~$\pm$~9.00 & 77.84~$\pm$~1.91 & 89.11~$\pm$~3.48 & 81.31~$\pm$~2.61 \\
        $\epsilon_{rob}$ = 0.05 & 76.62~$\pm$~2.22 & 17.15~$\pm$~0.63 & 45.78~$\pm$~12.10 & 76.83~$\pm$~2.73 & 90.73~$\pm$~2.79 & 79.54~$\pm$~3.36 \\
        $\epsilon_{rob}$ = 0.10 & 65.99~$\pm$~0.61 & 17.12~$\pm$~0.63 & 59.68~$\pm$~13.60 & 65.97~$\pm$~0.67 & 99.74~$\pm$~0.30 & 66.02~$\pm$~0.68 \\
        \midrule
        \multicolumn{7}{c}{\textbf{\cifar}}\\
            & $Acc_{te}$ ($\modeldef$) & $Acc_{rob}$ ($\modelstd$) & $Acc_{rob}$ ($\modeldef$) & $\modelstol$: $Acc_{te}$ & $\modelstol$: $Fid$ & $\modelstol$: $Fid_{cor}$ \\
        \midrule
        Baseline ($\modelstd$) & 83.30~$\pm$~1.93 & - & - & - & - & - \\
        $\epsilon_{rob}$ = 0.01 & 80.85~$\pm$~0.26 & 23.82~$\pm$~1.82 & 66.82~$\pm$~0.57 & 79.30~$\pm$~0.77 & 83.94~$\pm$~0.85 & 88.14~$\pm$~0.42 \\
        $\epsilon_{rob}$ = 0.03 & 70.69~$\pm$~0.44 & 9.12~$\pm$~0.76 & 38.85~$\pm$~0.84 & 70.31~$\pm$~1.21 & 79.17~$\pm$~0.73 & 79.82~$\pm$~0.78 \\
        $\epsilon_{rob}$ = 0.05 & 66.31~$\pm$~0.78 & 7.96~$\pm$~1.23 & 26.86~$\pm$~0.58 & 66.18~$\pm$~0.80 & 78.94~$\pm$~1.09 & 74.78~$\pm$~0.76 \\
        $\epsilon_{rob}$ = 0.10 & 62.65~$\pm$~15.41 & 5.20~$\pm$~1.66 & 19.29~$\pm$~7.96 & 61.43~$\pm$~15.63 & 86.77~$\pm$~1.26 & 66.62~$\pm$~17.33 \\
        \bottomrule

        \toprule
    \end{tabular}
\end{table*}


% Adversarial Training and Model Extraction
\noindent\textbf{\ref{advtr} (Adversarial Training) and~\ref{modelext} (Unauthorized Model Ownership).}
We evaluate how adversarial training (\ref{advtr}) influences unauthorized model ownership (\ref{modelext}), summarized in Table~\ref{tab:advtr_modelext}.
In this setting, the adversary extracts $\modelstol$ from a robustly trained model $\modeldef$.
We report results across four datasets and observe that the interaction varies by dataset.

On \census, adversarial training has little effect on unauthorized model ownership.
Across $\epsilon_{rob}=0.01$--$0.10$, $\modelstol$ fidelity increases slightly from $\approx95\%$ to $\approx98\%$, while correct fidelity remains stable, showing that the stolen model continues to mirror both the predictions and correctness of $\modeldef$.

For \fmnist, the trend is similar but slightly more pronounced.
Fidelity rises to $\approx96\%$, while correct fidelity declines from $\approx88\%$ to $\approx82\%$ as robustness increases.
Because $\modeldef$ continues to retain high test accuracy ($Acc_{te}$), $\modelstol$ remains a strong approximation, though its agreement on correct predictions weakens.

On more complex datasets, the pattern shifts.
For \lfw, $Fid$ rises from 88\% to nearly 100\%, but $Fid_{cor}$ drops from 84\% to 66\% as $\modeldef$ loses accuracy.
On \cifar, the effect is clearer: at $\epsilon_{rob}=0.10$, $\modeldef$ has $Acc_{te}\approx62\%$ and $Acc_{rob}\approx19\%$, while $\modelstol$ attains similar $Acc_{te}$ ($\approx61\%$) and $Fid$ ($\approx87\%$), roughly $8$ p.p. more than when $\epsilon_{rob}=0.03$.
Thus, the stolen model faithfully reproduces a boundary that is increasingly inaccurate.

Overall, adversarial training often improves raw fidelity but reduces correct fidelity, especially on harder datasets.
This suggests it may \emph{increase} susceptibility to unauthorized model ownership by producing more predictable decision boundaries.
While it strengthens robustness, it can also raise this susceptibility.
Further study with stronger defenses and modern architectures is needed to determine whether this effect is fundamental or an artifact of weaker training.

% Outlier Removal and Model Extraction
\begin{table*}[htbp]
\footnotesize
\caption{Interaction between~\ref{outrem} Outlier Removal and~\ref{modelext} Unauthorized Model Ownership. Removing outliers lowers $\modeldef$ accuracy but leaves stolen-model fidelity essentially unchanged.}
\label{tab:outrem_modext}
\begin{center}
\setlength\tabcolsep{2pt}
\begin{tabular}{ l c c c c c }
    \bottomrule

    \toprule
    & \multirow{2}{*}{$\modelstd$} & \multicolumn{4}{c}{$\modeldef$ (\% of outliers removed followed by retraining)} \\
        & & \textbf{10\%} & \textbf{20\%} & \textbf{30\%} & \textbf{40\%} \\
    \midrule
    \multicolumn{6}{c}{\textbf{\census}} \\
    $Acc_{te}$ & 84.60~$\pm$~0.16 & 81.80~$\pm$~0.30 & 80.73~$\pm$~0.52 & 79.43~$\pm$~0.45 & 77.61~$\pm$~0.39\\
    $\modelstol$: $Acc_{te}$ & 84.56~$\pm$~0.14 & 81.71~$\pm$~0.37 & 80.63~$\pm$~0.61 & 79.41~$\pm$~0.43 & 77.55~$\pm$~0.47\\
    $\modelstol$: $Fid$ & 99.12~$\pm$~0.08 & 98.88~$\pm$~0.24 & 99.07~$\pm$~0.14 & 99.17~$\pm$~0.19 & 99.00~$\pm$~0.19\\
    $\modelstol$: $Fid_{cor}$ & 84.14~$\pm$~0.17 & 81.19~$\pm$~0.36 & 80.22~$\pm$~0.62 & 79.00~$\pm$~0.43 & 77.08~$\pm$~0.37\\
    \midrule

    \multicolumn{6}{c}{\textbf{\lfw}}\\
    $Acc_{te}$ & 81.40~$\pm$~0.83  & 80.78~$\pm$~1.19 & 79.41~$\pm$~1.65 & 77.61~$\pm$~2.78 & 78.25~$\pm$~0.29\\
    $\modelstol$: $Acc_{te}$ & 80.96~$\pm$~0.48 & 80.08~$\pm$~1.16 & 78.77~$\pm$~2.24 & 78.03~$\pm$~2.53 & 76.09~$\pm$~1.04\\
    $\modelstol$: $Fid$ & 94.55~$\pm$~0.67 & 94.31~$\pm$~1.53 & 93.35~$\pm$~1.56 & 95.47~$\pm$~0.59 & 93.46~$\pm$~2.37 \\
    $\modelstol$: $Fid_{cor}$ & 78.46~$\pm$~0.38 & 77.58~$\pm$~0.66 & 75.76~$\pm$~1.96 & 75.56~$\pm$~2.56 & 73.90~$\pm$~1.64\\
    \midrule


    \multicolumn{6}{c}{\textbf{\fmnist}}\\
    $Acc_{te}$ & 83.51~$\pm$~0.55 & 79.19~$\pm$~0.58 & 78.40~$\pm$~1.09 & 75.39~$\pm$~1.22 & 75.22~$\pm$~1.27\\
    $\modelstol$: $Acc_{te}$ & 82.74~$\pm$~0.66 & 79.61~$\pm$~0.67 & 79.28~$\pm$~0.45 & 74.79~$\pm$~1.46 & 75.65~$\pm$~1.58\\
    $\modelstol$: $Fid$ & 92.89~$\pm$~0.45  & 90.61~$\pm$~0.34 & 89.68~$\pm$~0.56 & 89.72~$\pm$~0.59 & 88.75~$\pm$~0.85\\
    $\modelstol$: $Fid_{cor}$ & 80.27~$\pm$~0.58 & 75.65~$\pm$~0.70 & 74.66~$\pm$~0.90 & 71.03~$\pm$~1.48 & 70.91~$\pm$~1.62\\
    \midrule

    \multicolumn{6}{c}{\textbf{\cifar}}\\
    $Acc_{te}$ &  83.65~$\pm$~0.38 & 75.35~$\pm$~0.33 & 73.86~$\pm$~1.32 & 71.86~$\pm$~0.67 & 70.50~$\pm$~1.01\\
    $\modelstol$: $Acc_{te}$ & 82.03~$\pm$~0.56 & 75.81~$\pm$~0.48 & 74.26~$\pm$~1.32 & 72.13~$\pm$~1.11 & 71.62~$\pm$~1.25\\
    $\modelstol$: $Fid$ & 86.52~$\pm$~0.32 & 84.66~$\pm$~0.76 & 84.14~$\pm$~1.37 & 83.83~$\pm$~0.49 & 83.61~$\pm$~0.89\\
    $\modelstol$: $Fid_{cor}$ & 77.73~$\pm$~0.54 & 70.17~$\pm$~0.41 & 68.54~$\pm$~1.68 & 66.49~$\pm$~0.92 & 65.43~$\pm$~1.26\\
    \bottomrule

    \toprule
\end{tabular}
\end{center}
\end{table*}


\noindent\textbf{\ref{outrem} (Outlier Removal) and~\ref{modelext} (Unauthorized Model Ownership).}
Finally, we evaluate the interaction between outlier removal (\ref{outrem}) and unauthorized model ownership (\ref{modelext}).

\begin{figure}[!ht]
    \centering
    \includegraphics[width=\linewidth]{figures/fig_outrem_fid.pdf}
    \caption{Comparison of $Acc_{te}$ of $\modeldef$ and $Fid$ between $\modeldef$ and $\modelstol$ for the given percentage of outliers removed from $\modeldef$.
    $Acc_{te}$ shows a clear downward trend, while fidelity is not affected.
    The colors differentiate the dataset, while the dotted lines represent $Acc_{te}$ and the solid lines represent $Fid$. }
    \label{fig:outrem_fid}
    \Description{Line plot of defended-model test accuracy and stolen-model fidelity against the percentage of outliers removed, for three datasets. Test accuracy trends downward as more outliers are removed, while fidelity stays roughly constant.}
\end{figure}

Outlier removal trains the model after discarding training records with the lowest influence, estimated using $k$-nearest-neighbor Shapley values~\cite{jia2019EfficientTaskspecificData}.
These scores quantify how much each training record contributes to model predictions on a held-out set, allowing us to iteratively remove the least influential records before retraining the model on the remaining data.

Our results show that as the percentage of outliers removed increases, the test accuracy of the defended model ($\modeldef$) decreases.
This is expected since~\ref{outrem} uses Shapley values to identify data records with the highest influence, and removing them can degrade accuracy~\cite{jia2019EfficientTaskspecificData}.
For example, on \census, $Acc_{te}$ drops from $\approx84.6\%$ (baseline) to $\approx81\%$ with $10\%$ outliers removed, and further to $\approx77\%$ at $40\%$.

\begin{figure}[ht]
    \centering
    \includegraphics[width=\linewidth]{figures/fig_outrem_cor_fid.pdf}
    \caption{Comparison of $Acc_{te}$ and $Fid_{cor}$ for the given percentage of outliers removed from $\modeldef$.
    Both measures follow a similar downward trend.
    The colors differentiate the dataset, while the dotted lines represent $Acc_{te}$ and the solid lines represent $Fid_{cor}$. Error bars omitted for readability.}
    \label{fig:outrem_cor_fid}
    \Description{Line plot of defended-model test accuracy and correct-label fidelity against the percentage of outliers removed, for three datasets. Both measures follow a similar downward trend as more outliers are removed.}
\end{figure}

Analyzing the effect this has on unauthorized model ownership, Figure~\ref{fig:outrem_fid} shows that while $Acc_{te}$ of $\modeldef$ drops significantly for all datasets, the fidelity of $\modelstol$ stays constant (in the case of \census) or drops slightly (in the case of \cifar and \fmnist).
This suggests that outlier removal does not affect unauthorized model ownership.

Figure~\ref{fig:outrem_cor_fid} compares $Fid_{cor}$ with $Acc_{te}$, revealing that both trend downward.
For \cifar, removing $10\%$ of the outliers causes a $\geq7.5$~p.p. drop in both $Acc_{te}$ and $Fid_{cor}$.
This indicates that while overall fidelity is preserved, correct fidelity is sensitive to the degradation of $\modeldef${'s} accuracy.
We attribute this to the drop in $Acc_{te}$.
We conclude that~\ref{outrem} does not meaningfully affect~\ref{modelext}.

\begin{takeaway}
Each interaction here changes with the dataset and the defense strength, so it cannot be predicted from the algorithms alone.
It has to be measured.
\end{takeaway}
